% !TEX root = ../matan.tex

\section{Неравенство Минковского}

Пусть $\langle a,b\rangle\subset\overline{\RR}=[-\infty;+\infty]$ --- промежуток,
а $f,g$ интегрируемы по Риману на каждом $[c,d]\subset\langle a,b\rangle$. Для
$p\ge1$ положим
\[
\|f\|_{L^p}:=\left(\int_a^b |f(x)|^p\,dx\right)^{\frac1p}\in[0;+\infty].
\]
Неравенство Минковского --- это обобщённое неравенство треугольника для нормы
$\|\cdot\|_{L^p}$.

\begin{Theorem}{Неравенство Минковского}{}
	При $p\in[1;+\infty)$ выполнено
	\[
	\|f+g\|_{L^p}\le \|f\|_{L^p}+\|g\|_{L^p}.
	\]
\end{Theorem}

\begin{proof}
	\textbf{Случай $p=1$.} В силу $|f+g|\le|f|+|g|$ и линейности интеграла
	\[
	\int_a^b|f+g|\,dx\le\int_a^b\bigl(|f|+|g|\bigr)\,dx=\int_a^b|f|\,dx+\int_a^b|g|\,dx,
	\]
	то есть $\|f+g\|_{L^1}\le\|f\|_{L^1}+\|g\|_{L^1}$.

	\medskip
	\textbf{Случай $p\in(1;+\infty)$.} Если $\|f\|_{L^p}=+\infty$ или
	$\|g\|_{L^p}=+\infty$, правая часть бесконечна и неравенство тривиально.
	Считаем поэтому обе нормы конечными.

	\emph{Шаг 1: конечность $\|f+g\|_{L^p}$.} Докажем вспомогательное неравенство
	\[
	(s+t)^p\le 2^{p-1}(s^p+t^p),\qquad s,t\ge0.
	\]
	Функция $u\mapsto u^p$ при $p>1$ выпукла вниз на $[0;+\infty)$, поэтому
	$\left(\frac{s+t}{2}\right)^p\le\frac{s^p+t^p}{2}$; умножая на $2^p$, получаем
	требуемое. Применяя его к $s=|f(x)|$, $t=|g(x)|$ и интегрируя, имеем
	\[
	\int_a^b|f+g|^p\,dx\le\int_a^b\bigl(|f|+|g|\bigr)^p\,dx
	\le 2^{p-1}\left(\int_a^b|f|^p+\int_a^b|g|^p\right)<+\infty,
	\]
	то есть $\|f+g\|_{L^p}<+\infty$.

	\emph{Шаг 2: основная оценка.} Если $\|f+g\|_{L^p}=0$, доказывать нечего.
	Пусть $\|f+g\|_{L^p}>0$. Пусть $q$ сопряжён к $p$, то есть
	$\frac1p+\frac1q=1$, откуда $q=\frac{p}{p-1}$ и $(p-1)q=p$. Запишем
	\[
	\|f+g\|_{L^p}^p=\int_a^b|f+g|^p\,dx=\int_a^b|f+g|^{\,p-1}\,|f+g|\,dx
	\le\int_a^b|f+g|^{\,p-1}\bigl(|f|+|g|\bigr)\,dx,
	\]
	где использовано $|f+g|\le|f|+|g|$. Раскрывая скобки и применяя к каждому
	слагаемому неравенство Гельдера с показателями $p$ и $q$, получаем для
	первого интеграла
	\[
	\int_a^b|f|\,|f+g|^{\,p-1}\,dx
	\le\left(\int_a^b|f|^p\right)^{\frac1p}
	\left(\int_a^b|f+g|^{(p-1)q}\right)^{\frac1q}
	=\|f\|_{L^p}\,\|f+g\|_{L^p}^{\,p-1},
	\]
	так как $(p-1)q=p$ и $\frac1q=\frac{p-1}{p}$. Аналогично
	\[
	\int_a^b|g|\,|f+g|^{\,p-1}\,dx\le\|g\|_{L^p}\,\|f+g\|_{L^p}^{\,p-1}.
	\]
	Складывая, получаем
	\[
	\|f+g\|_{L^p}^p\le \|f+g\|_{L^p}^{\,p-1}\bigl(\|f\|_{L^p}+\|g\|_{L^p}\bigr).
	\]
	Поскольку $0<\|f+g\|_{L^p}<+\infty$, делим на $\|f+g\|_{L^p}^{\,p-1}$ и
	приходим к
	\[
	\|f+g\|_{L^p}\le \|f\|_{L^p}+\|g\|_{L^p}. \qedhere
	\]
\end{proof}

\begin{Remark}{}{}
	По индукции неравенство Минковского распространяется на любое конечное число
	слагаемых: $\bigl\|\sum_{i} f_i\bigr\|_{L^p}\le\sum_{i}\|f_i\|_{L^p}$.
\end{Remark}

\begin{Proposition}{Обратное неравенство Минковского}{}
	Если дополнительно $f(x)\ge0$ и $g(x)\ge0$ для всех $x$, то при $0<p\le1$
	\[
	\|f+g\|_{L^p}\ge \|f\|_{L^p}+\|g\|_{L^p}.
	\]
\end{Proposition}

\begin{proof}
	Доказательство повторяет схему прямого неравенства (с обратными
	неравенствами Юнга и Гельдера при $0<p<1$). Существенна неотрицательность
	$f,g$: без неё утверждение неверно. Контрпример: при $f\ne0$, $g=-f$ левая
	часть равна $0$, а правая равна $2\|f\|_{L^p}>0$, что противоречит
	неравенству.
\end{proof}

\begin{Proposition}{Неравенство Минковского для сумм}{}
	Пусть $a_k,b_k\ge0$, $k=1,\dots,N$, и $p\in[1;+\infty)$. Тогда
	\[
	\left(\sum_{k=1}^N (a_k+b_k)^p\right)^{\frac1p}
	\le
	\left(\sum_{k=1}^N a_k^p\right)^{\frac1p}
	+\left(\sum_{k=1}^N b_k^p\right)^{\frac1p}.
	\]
\end{Proposition}

\begin{proof}
	Сумма --- интеграл от ступенчатой функции. Возьмём $f,g$, постоянные на
	$[k-1,k]$ со значениями $a_k,b_k$ соответственно ($k=1,\dots,N$). Тогда
	\[
	\left(\int_0^N|f|^p\right)^{\frac1p}=\left(\sum a_k^p\right)^{\frac1p},
	\quad
	\left(\int_0^N|g|^p\right)^{\frac1p}=\left(\sum b_k^p\right)^{\frac1p},
	\quad
	\left(\int_0^N|f+g|^p\right)^{\frac1p}=\left(\sum (a_k+b_k)^p\right)^{\frac1p},
	\]
	и интегральное неравенство Минковского даёт требуемое.
\end{proof}
